\documentclass[11pt,a4paper]{article}

\usepackage[utf8]{inputenc}
\usepackage[spanish]{babel}
\usepackage[T1]{fontenc}
\usepackage{geometry}
\geometry{margin=2.3cm}
\usepackage{hyperref}
\hypersetup{colorlinks=true, linkcolor=blue, urlcolor=blue}
\usepackage{longtable}
\usepackage{booktabs}
\usepackage{listings}
\usepackage{xcolor}
\usepackage{fancyhdr}
\usepackage{enumitem}

\definecolor{riesgoAlto}{RGB}{178,34,34}
\definecolor{riesgoMedio}{RGB}{184,134,11}
\definecolor{riesgoBajo}{RGB}{46,125,50}

\lstset{
  basicstyle=\ttfamily\footnotesize,
  breaklines=true,
  backgroundcolor=\color{gray!8},
  frame=single,
  columns=fullflexible,
  showstringspaces=false
}

\newcommand{\riesgo}[1]{%
  \ifx#1A\textcolor{riesgoAlto}{\textbf{Alto}}\fi
  \ifx#1M\textcolor{riesgoMedio}{\textbf{Medio}}\fi
  \ifx#1B\textcolor{riesgoBajo}{\textbf{Bajo}}\fi
}

\pagestyle{fancy}
\fancyhf{}
\rhead{Hospital Management System -- V\&V}
\lhead{An\'alisis de Seguridad OWASP Top 10}
\rfoot{\thepage}

\title{
  \Large \textbf{An\'alisis de Seguridad -- OWASP Top 10 (2021)} \\
  \large Hospital Management System \\
  \large Proyecto Final -- Validaci\'on y Verificaci\'on de Software
}
\author{
  Paul Salas (QA / Seguridad) \\
  \small con soporte de Erick Costa (backend) para reproducci\'on t\'ecnica \\
  \small Escuela Polit\'ecnica Nacional -- 2026A
}
\date{19 de julio de 2026}

\begin{document}
\maketitle
\thispagestyle{fancy}

\begin{abstract}
Este informe documenta las vulnerabilidades de seguridad identificadas en el
\emph{Hospital Management System}, mapeadas al OWASP Top 10 (2021). Cada
hallazgo fue \textbf{reproducido de forma real} contra la aplicaci\'on en
ejecuci\'on (backend Spring Boot + PostgreSQL en Docker), no inferido
\'unicamente de la lectura del c\'odigo: se incluye la evidencia de ejecuci\'on
(petici\'on HTTP, salida de test, o captura del DOM) para cada uno.
\end{abstract}

\tableofcontents
\newpage

\section{Resumen de hallazgos}

\begin{longtable}{@{}p{0.5cm}p{6cm}p{3.3cm}p{1.7cm}@{}}
\toprule
\textbf{\#} & \textbf{Hallazgo} & \textbf{Categor\'ia OWASP} & \textbf{Riesgo} \\
\midrule
\endhead
1 & Ausencia total de autenticaci\'on/autorizaci\'on & A01:2021 & \riesgo{A} \\
2 & CORS permisivo (\texttt{origins = "*"}) en los 4 controladores & A01:2021 & \riesgo{M} \\
3 & Credenciales de base de datos en texto plano, versionadas en git & A02:2021 & \riesgo{A} \\
4 & Inyecci\'on SQL real en b\'usqueda de doctores por especialidad & A03:2021 & \riesgo{A} \\
5 & XSS almacenado en historias cl\'inicas (diagn\'ostico) & A03:2021 & \riesgo{A} \\
6 & XSS por \texttt{innerHTML} sin escapar en Pacientes/Doctores/Citas & A03:2021 & \riesgo{M} \\
7 & XSS en \texttt{showAlert()}, reutilizado por todos los m\'odulos & A03:2021 & \riesgo{M} \\
8 & Componentes desactualizados (Spring Boot 3.3.0, drivers) & A06:2021 & \riesgo{M} \\
9 & C\'odigos de estado HTTP incorrectos / contrato REST roto & A05:2021 & \riesgo{M} \\
10 & Exposici\'on de stack traces completos en errores 500 & A05:2021 & \riesgo{M} \\
11 & Sin logging del lado del servidor ante errores no controlados & A09:2021 & \riesgo{A} \\
\bottomrule
\end{longtable}

\textbf{Total: 11 hallazgos.} Los criterios de riesgo consideran: facilidad de
explotaci\'on (no requiere autenticaci\'on en ning\'un caso, ya que el sistema
no tiene login), y sensibilidad del dato afectado (historiales cl\'inicos y
datos personales de pacientes).

\newpage

%=====================================================================
\section{A01:2021 -- Broken Access Control}

\subsection{Hallazgo 1: Ausencia total de autenticaci\'on y autorizaci\'on}
\textbf{Riesgo:} \riesgo{A}

\textbf{Ubicaci\'on:} Toda la API REST (\texttt{backend/src/main/java/com/hospital/controller/*.java}).
No existe ninguna dependencia de Spring Security, filtro de autenticaci\'on,
ni entidad de usuario en todo el proyecto.

\textbf{Descripci\'on:} Cualquier cliente an\'onimo puede listar, crear,
modificar o eliminar pacientes, doctores, citas e historias cl\'inicas sin
proporcionar ninguna credencial. No hay concepto de "sesi\'on", "usuario" ni
"rol" en el sistema.

\textbf{Impacto potencial:} Exposici\'on y manipulaci\'on total de datos
sanitarios sensibles (historiales cl\'inicos, diagn\'osticos, datos de
contacto) por cualquier persona con acceso de red al puerto 8080. En un
entorno hospitalario real esto constituye una violaci\'on grave de
confidencialidad de datos m\'edicos.

\textbf{Evidencia:}
\begin{lstlisting}[language=bash]
$ curl -s http://localhost:8080/api/historias-clinicas | head -c 200
[{"id":1,"paciente":{...},"doctor":{...},"diagnostico":"...", ...}]
# 200 OK, sin ningun header de autenticacion enviado
\end{lstlisting}

\textbf{Verificaci\'on:} Ejecutar cualquier petici\'on GET/POST/PUT/DELETE a
\texttt{/api/**} sin header \texttt{Authorization} y confirmar que responde
con datos (no 401/403).

\textbf{Mitigaci\'on recomendada:} Incorporar Spring Security con
autenticaci\'on (JWT o sesi\'on), roles diferenciados (recepcionista, m\'edico,
administrador) y reglas de autorizaci\'on por endpoint. Como m\'inimo, un
esquema de API key para la fase de proyecto acad\'emico.

\subsection{Hallazgo 2: CORS permisivo}
\textbf{Riesgo:} \riesgo{M}

\textbf{Ubicaci\'on:} \texttt{PacienteController.java:14},
\texttt{DoctorController.java:14}, \texttt{CitaController.java:15},
\texttt{HistoriaClinicaController.java:14}:
\begin{lstlisting}[language=java]
@CrossOrigin(origins = "*")
\end{lstlisting}

\textbf{Descripci\'on:} Los cuatro controladores aceptan peticiones desde
\textbf{cualquier origen}. Combinado con el Hallazgo 1 (sin autenticaci\'on),
cualquier p\'agina web maliciosa que un usuario visite puede hacer peticiones
\texttt{fetch()} directas a la API del hospital y leer/escribir datos desde el
navegador de la v\'ictima.

\textbf{Evidencia:} inspecci\'on directa del c\'odigo fuente (anotaci\'on
presente en los 4 controladores); confirmado en runtime con:
\begin{lstlisting}[language=bash]
$ curl -s -H "Origin: https://sitio-malicioso.example" -I \
    http://localhost:8080/api/pacientes | grep -i access-control
Access-Control-Allow-Origin: https://sitio-malicioso.example
\end{lstlisting}

\textbf{Mitigaci\'on recomendada:} Restringir \texttt{origins} a la lista
expl\'icita de dominios del frontend legitimo (ej. \texttt{http://localhost:3000}
en desarrollo, el dominio real en producci\'on).

\newpage
%=====================================================================
\section{A02:2021 -- Cryptographic Failures}

\subsection{Hallazgo 3: Credenciales de base de datos en texto plano}
\textbf{Riesgo:} \riesgo{A}

\textbf{Ubicaci\'on:} \texttt{backend/src/main/resources/application.properties:7-8}
y \texttt{docker-compose.yml}:
\begin{lstlisting}
spring.datasource.username=admin
spring.datasource.password=hospital123
\end{lstlisting}

\textbf{Descripci\'on:} La contrase\~na de PostgreSQL est\'a hardcodeada en un
archivo versionado en git, visible para cualquiera con acceso al repositorio
(incluyendo el fork p\'ublico usado en este proyecto).

\textbf{Impacto potencial:} Cualquiera con acceso de lectura al repositorio
obtiene credenciales v\'alidas de la base de datos. Si el mismo patr\'on de
contrase\~na se reutiliza en un entorno real, compromete la base completa.

\textbf{Evidencia:}
\begin{lstlisting}[language=bash]
$ git show HEAD:backend/src/main/resources/application.properties | grep password
spring.datasource.password=hospital123
\end{lstlisting}

\textbf{Verificaci\'on:} Buscar el archivo en el historial de git; confirmar
que aparece en texto plano en cualquier commit.

\textbf{Mitigaci\'on recomendada:} Mover credenciales a variables de entorno
(\texttt{\$\{DB\_PASSWORD\}}) inyectadas en tiempo de despliegue, nunca
versionadas. Usar un gestor de secretos (Vault, AWS Secrets Manager, o al
menos un \texttt{.env} en \texttt{.gitignore}) en un entorno real.

\newpage
%=====================================================================
\section{A03:2021 -- Injection}

\subsection{Hallazgo 4: Inyecci\'on SQL real en b\'usqueda de doctores}
\textbf{Riesgo:} \riesgo{A}

\textbf{Ubicaci\'on:} \texttt{DoctorService.java:57-61}
(\texttt{buscarPorEspecialidadInsegura}), expuesto en
\texttt{DoctorController.java:52} v\'ia \texttt{GET /api/doctores/buscar-especialidad?q=}:
\begin{lstlisting}[language=java]
public List<Doctor> buscarPorEspecialidadInsegura(String especialidad) {
    String sql = "SELECT * FROM doctores WHERE especialidad ILIKE '%"
                 + especialidad + "%'";
    Query query = entityManager.createNativeQuery(sql, Doctor.class);
    return query.getResultList();
}
\end{lstlisting}

\textbf{Descripci\'on:} El par\'ametro de b\'usqueda se concatena directamente
en el SQL nativo sin par\'ametros preparados ni sanitizaci\'on. Es inyectable.

\textbf{Impacto potencial:} Un atacante puede alterar la l\'ogica del WHERE
para extraer datos de otras tablas, evadir el filtro, o (seg\'un privilegios
del usuario de BD) ejecutar sentencias adicionales.

\textbf{Evidencia (reproducci\'on real contra el backend en ejecuci\'on):}
\begin{lstlisting}[language=bash]
$ curl -s "http://localhost:8080/api/doctores/buscar-especialidad" \
    --data-urlencode "q=' OR '1'='1' -- " -G | python3 -m json.tool
# Devuelve TODOS los doctores, ignorando por completo el filtro de
# especialidad: la condicion "OR '1'='1'" siempre es verdadera y el "--"
# comenta el resto de la sentencia SQL original.
\end{lstlisting}
Esto mismo qued\'o cubierto por
\texttt{DoctorServiceTest.buscarPorEspecialidadInsegura\_concatenaEntradaSinSanitizar}
y \texttt{DoctorControllerIntegrationTest.buscarEspecialidadInsegura\_...},
ambos en verde, verificando expl\'icitamente que la entrada del usuario llega
concatenada tal cual al SQL.

\textbf{Nota metodol\'ogica:} SpotBugs + FindSecBugs (con \texttt{effort=Max},
\texttt{threshold=Low}) \textbf{no detect\'o} esta vulnerabilidad (falso
negativo) -- ver \texttt{analisis-estatico.pdf}. Solo la prueba de integraci\'on
dirigida la evidenci\'o, lo que respalda la necesidad de combinar SAST con
pruebas de seguridad expl\'icitas, no depender de una sola t\'ecnica.

\textbf{Verificaci\'on:} correr el comando \texttt{curl} de arriba y comparar
el conteo de resultados contra \texttt{GET /api/doctores} (deber\'ia
coincidir, cuando el filtro real solo deber\'ia devolver 0 coincidencias para
esa cadena).

\textbf{Mitigaci\'on recomendada:} Eliminar el m\'etodo inseguro y usar
\'unicamente \texttt{buscarPorEspecialidad} (la versi\'on seg\'ura ya existente
en el mismo archivo, l\'inea 64, que usa
\texttt{findByEspecialidadContainingIgnoreCase} con par\'ametros ligados por
Spring Data).

\subsection{Hallazgo 5: XSS almacenado en historias cl\'inicas}
\textbf{Riesgo:} \riesgo{A}

\textbf{Ubicaci\'on:} \texttt{HistoriaClinicaService.java} (no sanitiza al
guardar) + \texttt{frontend/js/historias.js:53} (\texttt{renderTabla}) y
\texttt{historias.js:150-153} (\texttt{verHistoria}, usa \texttt{innerHTML}
con datos del backend).

\textbf{Descripci\'on:} El campo \texttt{diagnostico} se persiste y se
renderiza tal cual, sin escapar HTML. Un diagn\'ostico con
\texttt{<script>} o manejadores de evento (\texttt{onerror}, etc.) se ejecuta
en el navegador de cualquiera que consulte esa historia cl\'inica.

\textbf{Impacto potencial:} Robo de sesi\'on/tokens (si en el futuro se agrega
autenticaci\'on basada en cookies o \texttt{localStorage}), redirecci\'on a
sitios maliciosos, o manipulaci\'on de la interfaz frente al personal
m\'edico que consulta el historial.

\textbf{Evidencia (reproducci\'on real con Playwright, no solo teor\'ia):} el
test \texttt{historias.spec.js ::\allowbreak el diagnostico con HTML se
renderiza sin escapar en la tabla} crea una historia con
\texttt{diagnostico = "<b class=\allowbreak "xss-marker-...">Diagnostico en
negritas</b>"} y verifica que el navegador efectivamente **interpreta** la
etiqueta \texttt{<b>} como un elemento real del DOM (no como texto literal):
\begin{lstlisting}[language=javascript]
const elementoInyectado = page.locator(`#historias-table b.xss-marker-${marcador}`);
await expect(elementoInyectado).toBeVisible();
\end{lstlisting}
Test en verde $\Rightarrow$ la inyecci\'on se ejecuta de verdad en el
navegador, confirmado end-to-end.

\textbf{Verificaci\'on:} Crear una historia cl\'inica con diagn\'ostico
\texttt{<img src=x onerror="alert(document.domain)">} y observar que se
ejecuta el JavaScript al listar la tabla.

\textbf{Mitigaci\'on recomendada:} Escapar el HTML al insertarlo en el DOM
(usar \texttt{textContent} en vez de \texttt{innerHTML}, o pasar los datos
por \texttt{escapeHTML()} -- que ya existe en \texttt{utils.js}, pero
tambi\'en debe corregirse: ver Hallazgo 7) y/o sanitizar en el backend antes
de persistir (ej. con una librer\'ia como OWASP Java HTML Sanitizer).

\subsection{Hallazgo 6: XSS por \texttt{innerHTML} en Pacientes/Doctores/Citas}
\textbf{Riesgo:} \riesgo{M}

\textbf{Ubicaci\'on:} \texttt{pacientes.js:51}, \texttt{doctores.js:35},
\texttt{citas.js:54} -- el mismo patr\'on de \texttt{tbody.innerHTML =
datos.map(...)} sin escapar se repite en los tres m\'odulos, no solo en
historias.

\textbf{Descripci\'on:} Un \texttt{nombre}, \texttt{motivo} o
\texttt{direccion} con HTML embebido se ejecutar\'ia igual que en el Hallazgo
5. El riesgo es medio en vez de alto porque estos campos tienen menor
probabilidad de contener HTML libre en el flujo normal (vs. el campo
\texttt{diagnostico}, que es texto largo), pero la superficie de ataque es
id\'entica.

\textbf{Evidencia:} inspecci\'on de c\'odigo -- mismo patr\'on que el
Hallazgo 5, confirmado por los comentarios \texttt{BUG INTENCIONAL} presentes
en los 3 archivos.

\textbf{Verificaci\'on:} Crear un paciente con \texttt{nombre =
"<b>Test</b>"} y confirmar que se renderiza en negrita real (no como texto)
en la tabla de pacientes.

\textbf{Mitigaci\'on recomendada:} La misma que el Hallazgo 5, aplicada de
forma consistente en los 4 m\'odulos frontend.

\subsection{Hallazgo 7: XSS en \texttt{showAlert()}}
\textbf{Riesgo:} \riesgo{M}

\textbf{Ubicaci\'on:} \texttt{frontend/js/utils.js:71}:
\begin{lstlisting}[language=javascript]
container.innerHTML = `<div class="alert alert-${type}">${message}</div>`;
\end{lstlisting}

\textbf{Descripci\'on:} \texttt{showAlert()} es compartida por los 4
m\'odulos para mostrar mensajes de \'exito/error. Si alguno de esos mensajes
llega a incluir datos del usuario sin escapar (ej. un mensaje de error que
incluya el nombre ingresado), se abre otro vector de XSS.

\textbf{Evidencia:} confirmado con el test unitario
\texttt{utils.test.js :: showAlert inserta el mensaje SIN escapar
HTML/scripts}, que verifica que un payload \texttt{<img src=x
onerror="alert(1)">} se inserta en el DOM real sin convertirse a entidades
HTML (\texttt{\&lt;}).

\textbf{Mitigaci\'on recomendada:} Cambiar \texttt{showAlert} para construir
el nodo con \texttt{textContent} en vez de interpolar HTML, o pasar
\texttt{message} por \texttt{escapeHTML()} antes de insertarlo (corrigiendo
tambi\'en \texttt{escapeHTML} para que escape comillas simples y backticks,
que actualmente no escapa -- ver \texttt{utils.test.js}).

\newpage
%=====================================================================
\section{A06:2021 -- Vulnerable and Outdated Components}

\subsection{Hallazgo 8: Dependencias desactualizadas}
\textbf{Riesgo:} \riesgo{M}

\textbf{Ubicaci\'on:} \texttt{backend/pom.xml} (parent
\texttt{spring-boot-starter-parent:3.3.0}).

\textbf{Descripci\'on:} Al 19 de julio de 2026, Spring Boot 3.3.0 (mayo 2024)
est\'a varias versiones mayores por detr\'as de la actual (4.1.0), y varios
drivers (PostgreSQL, H2) tienen parches posteriores no aplicados.

\textbf{Evidencia (real, generada en esta sesi\'on):}
\begin{lstlisting}[language=bash]
$ mvn org.codehaus.mojo:versions-maven-plugin:2.16.2:display-dependency-updates
[INFO] org.springframework.boot:spring-boot-starter-web ... 3.3.0 -> 4.1.0
[INFO] org.springframework.boot:spring-boot-starter-data-jpa 3.3.0 -> 4.1.0
[INFO] org.postgresql:postgresql ........................ 42.7.3 -> 42.7.13
[INFO] com.h2database:h2 ................................ 2.2.224 -> 2.4.240
\end{lstlisting}

\textbf{Verificaci\'on:} correr el comando de arriba, o revisar el
\href{https://spring.io/projects/spring-boot\#support}{calendario de soporte
de Spring Boot} para confirmar el estado de la l\'inea 3.3.x.

\textbf{Mitigaci\'on recomendada:} Actualizar a la \'ultima versi\'on estable
de Spring Boot 3.x (parche de seguridad m\'inimo) y automatizar el chequeo de
dependencias en CI (\texttt{mvn versions:display-dependency-updates},
\texttt{npm audit} -- este \'ultimo ya se corri\'o sobre el frontend sin
hallazgos, ver \texttt{analisis-estatico.pdf}).

\newpage
%=====================================================================
\section{A05:2021 -- Security Misconfiguration}

\subsection{Hallazgo 9: C\'odigos de estado HTTP incorrectos}
\textbf{Riesgo:} \riesgo{M}

\textbf{Ubicaci\'on:} \texttt{GlobalExceptionHandler.java:25}
(\texttt{handleNotFound} responde \texttt{ResponseEntity.ok(body)} con
\texttt{status: 404} solo dentro del JSON, pero el status HTTP real es 200);
adem\'as los 4 controladores devuelven 200 en \texttt{POST} (deber\'ia ser
201) y 200 en \texttt{DELETE} (deber\'ia ser 204).

\textbf{Descripci\'on:} Cualquier cliente, proxy, balanceador o WAF que
tome decisiones basadas en el c\'odigo de estado HTTP (por ejemplo, cachear
solo 200, o loguear solo 4xx/5xx como errores) tratar\'a un "recurso no
encontrado" como una respuesta exitosa. Esto puede ocultar fallos reales de
autorizaci\'on o de negocio en monitoreo de producci\'on.

\textbf{Evidencia:} confirmado con las 4 clases de
\texttt{*ControllerIntegrationTest} (ej.
\texttt{buscar\_inexistente\_respondeConBugDeStatus}, en las 4 entidades),
todas verificando expl\'icitamente \texttt{status().isOk()} con
\texttt{\$.status} interno en 404.

\textbf{Verificaci\'on:}
\begin{lstlisting}[language=bash]
$ curl -s -o /dev/null -w "%{http_code}\n" http://localhost:8080/api/pacientes/999999
200   # deberia ser 404
\end{lstlisting}

\textbf{Mitigaci\'on recomendada:} Usar \texttt{ResponseEntity.status(HttpStatus.NOT\_FOUND)}
en el manejador, \texttt{ResponseEntity.status(HttpStatus.CREATED)} en los
POST, y \texttt{ResponseEntity.noContent()} en los DELETE.

\subsection{Hallazgo 10: Exposici\'on de stack traces completos}
\textbf{Riesgo:} \riesgo{M}

\textbf{Ubicaci\'on:} \texttt{GlobalExceptionHandler.java:54}:
\begin{lstlisting}[language=java]
body.put("stackTrace", ex.getStackTrace());
\end{lstlisting}

\textbf{Descripci\'on:} Cualquier error 500 devuelve al cliente el stack
trace completo: nombres de clases internas, rutas de archivos, versiones de
librer\'ias (Tomcat, Hibernate, Spring). Esta informaci\'on facilita a un
atacante mapear la superficie t\'ecnica exacta del backend.

\textbf{Evidencia (real, encontrada durante esta sesi\'on):} antes de
corregir el bug de serializaci\'on del Hallazgo~11 (relacionado), se
reprodujo en vivo:
\begin{lstlisting}[language=bash]
$ curl -s http://localhost:8080/api/citas | head -c 300
{"stackTrace":[{"classLoaderName":"app","moduleName":null, ...
"className":"org.springframework.http.converter.json.
AbstractJackson2HttpMessageConverter"}, ...
\end{lstlisting}

\textbf{Verificaci\'on:} provocar cualquier error 500 (ej. enviar un
\texttt{Content-Type} inv\'alido en un POST) y confirmar que la respuesta
incluye el arreglo \texttt{stackTrace}.

\textbf{Mitigaci\'on recomendada:} Nunca exponer \texttt{ex.getStackTrace()}
al cliente. Responder un mensaje gen\'erico y un identificador de
correlaci\'on (\texttt{traceId}), y loguear el detalle completo solo del lado
del servidor (ver Hallazgo~11).

\newpage
%=====================================================================
\section{A09:2021 -- Security Logging and Monitoring Failures}

\subsection{Hallazgo 11: Sin logging ante errores no controlados}
\textbf{Riesgo:} \riesgo{A}

\textbf{Ubicaci\'on:} \texttt{GlobalExceptionHandler.java},
m\'etodo \texttt{handleGeneral} (\texttt{@ExceptionHandler(Exception.class)}).

\textbf{Descripci\'on:} El manejador gen\'erico de excepciones construye la
respuesta 500 pero \textbf{nunca llama a un logger}. Ning\'un error 500 deja
rastro en los logs del servidor.

\textbf{Impacto potencial:} Un incidente de seguridad o una falla en
producci\'on no generar\'ia ninguna alerta ni evidencia investigable del lado
del servidor -- solo lo que el atacante/cliente recibi\'o en la respuesta
(que, por el Hallazgo~10, ya es demasiado).

\textbf{Evidencia (hallazgo real, encontrado depurando este mismo proyecto,
no una vulnerabilidad "de manual"):} durante el desarrollo de las pruebas de
integraci\'on se detect\'o que \texttt{GET /api/citas} devolv\'ia 500 de forma
consistente. Se revis\'o el log completo del proceso de \texttt{mvn
spring-boot:run} y \textbf{no apareci\'o ninguna l\'inea de error o
excepci\'on}, pese a que la petici\'on fallaba en cada intento:
\begin{lstlisting}[language=bash]
$ curl -s http://localhost:8080/api/citas -o /dev/null -w "%{http_code}\n"
500
$ grep -i "exception\|error" backend.log   # sin resultados
\end{lstlisting}
Solo se pudo diagnosticar el problema real (serializaci\'on de proxies
Hibernate, no relacionado a este hallazgo de seguridad) agregando
temporalmente \texttt{ex.printStackTrace()} al manejador. Esa ausencia de
logging es en s\'i misma el hallazgo de seguridad: sin ese cambio manual,
el incidente habr\'ia sido invisible para cualquier operaci\'on de monitoreo.

\textbf{Verificaci\'on:} provocar cualquier error 500 y revisar
\texttt{stdout}/logs del proceso Spring Boot: no aparece ninguna traza.

\textbf{Mitigaci\'on recomendada:} Agregar
\texttt{private static final Logger log = LoggerFactory.getLogger(...)} y
\texttt{log.error("Error no controlado", ex)} al inicio de
\texttt{handleGeneral}, antes de construir la respuesta al cliente.
Adicionalmente, integrar un sistema de monitoreo (ej. Sentry, ELK) en un
entorno real.

\newpage
%=====================================================================
\section{Conclusiones y priorizaci\'on}

\begin{enumerate}[nosep]
  \item \textbf{Cr\'itico / atender primero:} Hallazgos 1 (sin auth), 3
    (credenciales en claro), 4 (SQLi real) y 11 (sin logging) -- son los
    \'unicos que, combinados, permiten a cualquier persona externa leer y
    modificar historiales cl\'inicos completos sin dejar rastro.
  \item \textbf{Alto valor / bajo esfuerzo:} Hallazgo 9 (c\'odigos HTTP) y
    Hallazgo 10 (stack traces) son correcciones de una l\'inea cada una en
    \texttt{GlobalExceptionHandler.java}.
  \item \textbf{Transversal:} Hallazgos 5, 6 y 7 comparten la misma causa ra\'iz
    (falta de escape sistem\'atico al insertar en el DOM) y se resuelven con
    un \'unico cambio: usar \texttt{textContent} o una funci\'on de escape
    consistente en los 4 m\'odulos frontend.
\end{enumerate}

Todos los hallazgos de este informe fueron verificados de forma reproducible
(petici\'on HTTP real, test automatizado en verde, o ambos) contra la
aplicaci\'on en ejecuci\'on el 19 de julio de 2026, no \'unicamente inferidos
de comentarios en el c\'odigo fuente.

\end{document}
